% 图契约：三次测量重建三端网络；沿用可恢复原图的真实变量及结构，图号由正文图注统一给定。\n% 图契约：保留原稿的物理与电路拓扑，启用中文字体及中文断行，逐图检查标注。
\documentclass[tikz,border=14pt]{standalone}
\usepackage{fontspec,xeCJK}
\usepackage{amsmath,amssymb}
\usepackage{xcolor}
\usepackage{tikz}
\usetikzlibrary{arrows.meta,calc,patterns,positioning,decorations.pathmorphing,bending}
\IfFontExistsTF{Heiti SC}{\IfFontExistsTF{Heiti SC}{\setmainfont{Heiti SC}}{\setmainfont{FandolHei-Regular}}\IfFontExistsTF{Heiti SC}{\setsansfont{Heiti SC}}{\setsansfont{TeX Gyre Heros}}}{\IfFontExistsTF{Songti SC}{\setmainfont{Songti SC}}{\setmainfont{FandolSong-Regular}}\IfFontExistsTF{Songti SC}{\setsansfont{Songti SC}}{\setsansfont{FandolSong-Regular}}}
\definecolor{bluefield}{RGB}{24,76,255}
\definecolor{redfield}{RGB}{232,49,35}
\definecolor{orangefield}{RGB}{224,116,10}
\definecolor{greenfield}{RGB}{18,132,86}
\definecolor{purplefield}{RGB}{112,72,190}
\definecolor{linegray}{RGB}{60,68,78}
\definecolor{softblue}{RGB}{238,243,255}
\definecolor{softgreen}{RGB}{238,249,244}
\definecolor{softorange}{RGB}{255,247,235}
\definecolor{softpurple}{RGB}{247,242,255}
\tikzset{
  arrow/.style={-{Latex[length=3.1mm,width=2.1mm]},line width=1.1pt},
  thinarr/.style={-{Latex[length=2.5mm,width=1.7mm]},line width=.85pt},
  axis/.style={-{Latex[length=2.5mm,width=1.7mm]},line width=.8pt,draw=linegray},
  guide/.style={draw=linegray!55,dashed,line width=.65pt},
  fieldline/.style={-{Latex[length=2.5mm,width=1.7mm]},line width=1pt,draw=bluefield},
  box/.style={draw=linegray,rounded corners=3pt,fill=white,align=center,inner sep=6pt,line width=.75pt},
  bluebox/.style={box,draw=bluefield,fill=softblue},
  greenbox/.style={box,draw=greenfield,fill=softgreen},
  orangebox/.style={box,draw=orangefield,fill=softorange},
  purplebox/.style={box,draw=purplefield,fill=softpurple}
}
\IfFontExistsTF{Heiti SC}{\setCJKmainfont[AutoFakeBold=1.5]{Heiti SC}}{\setCJKmainfont[AutoFakeBold=1.5]{FandolHei-Regular}}
\IfFontExistsTF{Heiti SC}{\setCJKsansfont{Heiti SC}}{\setCJKsansfont{FandolHei-Regular}}
\begin{document}

\begin{tikzpicture}[x=1cm,y=1cm]
\path[use as bounding box] (-10,-4.6) rectangle (10,4.6);
\fill[white] (-10,-4.6) rectangle (10,4.6);
\node[font=\Large\bfseries,text width=18.4cm,align=center] at (0,4.05) {三端黑箱：三次悬空测量重建多端等效网络};
\node[bluebox,text width=4.2cm] at (-6.2,.3) {三端黑箱\\端点 $A,B,C$};
\node[greenbox,text width=5cm] at (0,1.6) {$C$ 悬空：测 $R_{AB}$\\$A$ 悬空：测 $R_{BC}$};
\node[orangebox,text width=5cm] at (0,-1.45) {$B$ 悬空：测 $R_{CA}$\\先重建星形臂};
\begin{scope}[shift={(6.4,.2)}]
  \coordinate (A) at (0,2.2);\coordinate (B) at (-2,-1.5);\coordinate (C) at (2,-1.5);
  \draw[linegray,line width=1pt] (A)--(B)--(C)--cycle;
  \node[above] at (A) {$A$};\node[below left] at (B) {$B$};\node[below right] at (C) {$C$};
  \node at (0,-2.5) {再换成等效三角形};
\end{scope}
\draw[arrow,bluefield] (-3.7,.3)--(-2.7,.8); \draw[arrow,greenfield] (2.8,.3)--(3.7,.3);
\end{tikzpicture}
\end{document}
